\chapter{Introdução} \label{Introducao}

Modelos de aprendizado de máquina fazem parte de um número crescente de aplicações cotidianas. Teclados de celulares sugerem a próxima palavra a ser digitada, plataformas de vídeo recomendam conteúdos com base no histórico de cada usuário e sistemas de apoio ao diagnóstico analisam exames de imagem em hospitais. Em comum, essas aplicações dependem de grandes volumes de dados para atingir bom desempenho, e esses dados frequentemente contêm informações pessoais e sensíveis, como mensagens privadas, fotografias e registros clínicos \cite{Goodfellow2016}.

No treinamento convencional, todos os dados são reunidos em um servidor central, onde o modelo é ajustado. Essa centralização traz dois problemas. O primeiro é o risco de vazamento: um único servidor que armazena dados de milhões de pessoas torna-se um alvo valioso para ataques. O segundo é de ordem legal e institucional, pois muitas organizações não podem, ou não desejam, transferir os dados de seus usuários para terceiros. Um hospital, por exemplo, dificilmente pode enviar prontuários de pacientes para ser processados fora de sua infraestrutura.

O aprendizado federado surgiu como uma alternativa a esse modelo centralizado \cite{McMahan2017}. Nele, diversos participantes, chamados de clientes, colaboram para treinar um modelo comum sem compartilhar seus dados brutos. Uma analogia útil é a de vários hospitais que desejam construir juntos um modelo de apoio ao diagnóstico. Em vez de enviar os prontuários de seus pacientes a uma central, cada hospital treina o modelo localmente com seus próprios dados e envia apenas as lições aprendidas, isto é, os ajustes que sugere aos parâmetros do modelo. Um servidor combina esses ajustes e devolve aos hospitais uma versão aprimorada do modelo, repetindo o processo por várias rodadas.

Os ajustes enviados pelos clientes são, em geral, gradientes. O gradiente indica em que direção cada parâmetro do modelo deve ser alterado para reduzir o erro nas amostras de treinamento do cliente. Na forma mais simples do aprendizado federado, o servidor atualiza o modelo combinando os gradientes recebidos da seguinte maneira:

\begin{equation} \label{eq:intro-fedsgd}
w_{t+1} = w_t - \eta \sum_{k=1}^{M} \frac{n_k}{n} \, g_k
\end{equation}

\noindent em que $w_t$ representa os parâmetros do modelo na rodada $t$, $\eta$ é a taxa de aprendizado, $M$ é o número de clientes, $n_k$ é a quantidade de amostras do cliente $k$, $n$ é o total de amostras e $g_k$ é o gradiente calculado pelo cliente $k$ sobre seus dados locais. A Equação~\ref{eq:intro-fedsgd} evidencia que o servidor nunca recebe as amostras em si, apenas os vetores $g_k$. A premissa implícita do aprendizado federado é que esses vetores contêm menos informação sobre os dados do que os próprios dados.

Essa premissa, contudo, mostrou-se frágil. Voltando à analogia, é como se, a partir apenas dos ajustes enviados por um hospital, fosse possível reconstruir os prontuários que os originaram. \citeonline{Zhu2019} demonstraram exatamente isso em um ataque denominado Deep Leakage from Gradients (DLG). O adversário cria uma imagem aleatória, calcula o gradiente que ela produziria no modelo e a ajusta repetidamente até que esse gradiente coincida com o gradiente compartilhado pelo cliente. Ao final do processo, a imagem artificial torna-se praticamente idêntica à imagem original de treinamento. Trabalhos posteriores tornaram esses ataques, conhecidos como ataques de inversão de gradientes, mais estáveis e eficazes. O iDLG mostrou que o rótulo de uma amostra pode ser extraído diretamente do gradiente, sem nenhuma otimização \cite{Zhao2020}, e outros métodos obtiveram reconstruções reconhecíveis de imagens de alta resolução e até de lotes com dezenas de imagens \cite{Geiping2020,Yin2021}. Compartilhar gradientes, portanto, não equivale a proteger os dados que os geraram.

Para mitigar esse risco, a literatura recorre principalmente a duas famílias de técnicas, que atuam modificando o gradiente antes que ele seja utilizado ou compartilhado. A primeira é a privacidade diferencial, cujo representante mais conhecido no treinamento de redes neurais é o DP-SGD \cite{Abadi2016}. Esse método limita a influência de cada amostra, recortando os gradientes que ultrapassam uma norma máxima, e em seguida adiciona ruído aleatório, de modo que a contribuição individual de cada dado fique mascarada. A segunda família é a compressão de gradientes, originalmente concebida para reduzir o custo de comunicação em treinamento distribuído. Nas esparsificações Top-$k$ \cite{Stich2018} e Random-$k$ \cite{Wangni2018}, apenas uma fração dos valores do gradiente é mantida, escolhida pelas maiores magnitudes ou por sorteio; na quantização QSGD \cite{Alistarh2017}, todos os valores são mantidos, mas arredondados para poucos níveis possíveis.

Ambas as famílias reduzem a informação disponível ao adversário, mas também alteram a direção em que o modelo é ajustado durante o treinamento. Quanto mais ruído é adicionado ou quanto mais valores são descartados, mais difícil se torna o ataque, porém menor tende a ser a acurácia do modelo final. Estabelece-se, assim, um compromisso entre privacidade e utilidade: uma defesa que protege completamente os dados, mas produz um modelo inútil, não resolve o problema, assim como um modelo excelente que expõe os dados de seus usuários também não é aceitável.

Apesar da diversidade de defesas propostas, comparações empíricas sistemáticas entre elas permanecem escassas. Os trabalhos que introduzem ataques avaliam defesas de forma incidental, geralmente com uma única arquitetura de rede \cite{Zhu2019}. Os trabalhos sobre privacidade diferencial quantificam a perda de acurácia, mas não a resistência empírica a ataques de inversão \cite{Abadi2016}. As técnicas de compressão, por sua vez, foram avaliadas quanto à economia de comunicação, sem análise de privacidade \cite{Stich2018,Wangni2018,Alistarh2017,Lin2018}. Mesmo as avaliações comparativas existentes concentram-se em arquiteturas convolucionais e em conjuntos específicos de defesas \cite{Wei2020,Huang2021}. Como cada estudo emprega bases de dados, arquiteturas e protocolos de ataque distintos, não é possível determinar, a partir da literatura, quais configurações oferecem o melhor equilíbrio entre proteção e desempenho, nem se as conclusões obtidas em um cenário se mantêm em outro.

Diante desse cenário, este trabalho busca responder à seguinte questão: sob um mesmo protocolo experimental, qual é o compromisso entre privacidade e utilidade proporcionado por mecanismos de privacidade diferencial e de compressão de gradientes contra ataques de inversão, e como esse compromisso varia de acordo com a base de dados, a arquitetura do modelo e o tamanho do lote de treinamento? A hipótese central é que mecanismos que reduzem a informação contida nos gradientes dificultam a reconstrução dos dados, elevando o erro de reconstrução do ataque, mas com impactos distintos sobre a acurácia do modelo, a depender da natureza da redução aplicada, seja ela a adição de ruído, o descarte de valores ou a perda de precisão numérica.

\section{Justificativa}

A relevância do tema decorre, em primeiro lugar, da crescente adoção do aprendizado federado em cenários que envolvem dados sensíveis, como saúde, finanças e dispositivos móveis \cite{Kairouz2021}. Nesses contextos, a garantia de que os dados não saem do dispositivo do usuário é frequentemente apresentada como suficiente para proteger sua privacidade. Os ataques de inversão de gradientes mostram que essa garantia é incompleta, e a escolha de uma defesa adequada torna-se uma decisão de engenharia com consequências diretas para as pessoas cujos dados são utilizados.

Em segundo lugar, há uma motivação regulatória. No Brasil, a Lei Geral de Proteção de Dados Pessoais estabelece que o tratamento de dados pessoais deve adotar medidas técnicas aptas a protegê-los de acessos não autorizados \cite{Brasil2018}. Compreender empiricamente a proteção oferecida por cada mecanismo de defesa, e o custo que ele impõe à qualidade do modelo, fornece subsídios objetivos para que desenvolvedores e organizações escolham técnicas compatíveis com esse tipo de exigência.

Em terceiro lugar, do ponto de vista científico, a ausência de um protocolo comum de avaliação dificulta a comparação entre defesas e pode levar a conclusões pouco generalizáveis. Uma defesa que parece eficaz em imagens simples de dígitos manuscritos, avaliada em uma única arquitetura, pode comportar-se de forma muito diferente em fotografias coloridas ou em redes baseadas em mecanismos de atenção. Um estudo que compare, sob as mesmas condições, defesas de famílias diferentes, em arquiteturas e bases de dados de complexidade distinta, contribui para preencher essa lacuna e para orientar pesquisas futuras, como a investigação de defesas combinadas.

Por fim, o tema reúne áreas centrais da formação em Ciência da Computação, como aprendizado de máquina, otimização, estatística e segurança da informação, e oferece a oportunidade de aplicar esses conhecimentos a um problema atual e de impacto social.

A escolha do tema também reflete o interesse pessoal da autora por segurança da informação e por sistemas distribuídos. O aprendizado federado situa-se exatamente na interseção dessas duas áreas: trata-se de um sistema distribuído, em que múltiplos participantes cooperam por meio da troca de mensagens com um servidor, e cuja principal promessa é de natureza de segurança, a de preservar a privacidade dos dados de cada participante. Investigar até que ponto essa promessa se sustenta diante de ataques reais, e quais mecanismos de defesa a tornam mais efetiva, permitiu unir esses dois interesses em um único problema de pesquisa.

\section{Objetivos}

O objetivo geral deste trabalho é avaliar empiricamente o compromisso entre privacidade e utilidade de mecanismos de defesa baseados em privacidade diferencial e em compressão de gradientes contra ataques de inversão de gradientes, sob um protocolo experimental uniforme que abrange diferentes bases de dados, arquiteturas de redes neurais e tamanhos de lote.

Para alcançar esse objetivo, o trabalho parte da implementação de cinco configurações de defesa aplicadas aos gradientes durante o treinamento: a ausência de defesa, que serve como referência, o DP-SGD, as esparsificações Top-$k$ e Random-$k$ e a quantização QSGD. Com essas defesas, busca-se treinar modelos de três arquiteturas com características distintas, o perceptron multicamadas (MLP), a rede neural convolucional (CNN) e o Vision Transformer (ViT), sobre três bases de imagens de dificuldade crescente, MNIST, Fashion-MNIST e CIFAR-10, variando o tamanho do lote e repetindo cada configuração diversas vezes para garantir a confiabilidade dos resultados.

Sobre os modelos treinados, pretende-se executar o ataque de inversão de gradientes iDLG e medir a privacidade de duas formas complementares: pelo erro quadrático médio entre as imagens originais e as reconstruídas, que indica o quanto da imagem foi recuperado, e pela taxa de acerto do rótulo inferido pelo ataque, que indica se a classe da amostra foi revelada. Paralelamente, a utilidade de cada modelo é medida pela acurácia no conjunto de teste, permitindo verificar o custo que cada defesa impõe ao desempenho.

Por fim, espera-se comparar estatisticamente cada defesa com o cenário sem defesa, utilizando testes de hipótese e medidas de tamanho de efeito, de modo a identificar as configurações que oferecem o melhor equilíbrio entre privacidade e utilidade e a discutir como esse equilíbrio depende da base de dados, da arquitetura e do tamanho do lote.

\section{Método do Trabalho}

O trabalho adota uma abordagem experimental e quantitativa, organizada em cinco etapas encadeadas. Cada etapa produz os insumos necessários para a seguinte, de forma que as conclusões finais possam ser rastreadas até as decisões tomadas no início do estudo.

A primeira etapa consistiu na revisão da literatura sobre aprendizado federado, ataques de inversão de gradientes e mecanismos de defesa. Esse estudo permitiu selecionar as técnicas avaliadas, representando as duas principais famílias de defesa, escolher o iDLG como ataque de referência, por sua simplicidade e eficácia no cenário de um único cliente, e definir as métricas utilizadas para medir privacidade e utilidade.

Na segunda etapa, foi planejado um experimento fatorial completo, no qual todas as combinações de quatro fatores são avaliadas: a arquitetura, com três opções; a base de dados, com três opções; a defesa, com cinco configurações; e o tamanho do lote, com lotes de uma e de trinta e duas amostras. Esse desenho resulta em 90 combinações, e cada uma delas foi executada com 30 sementes aleatórias distintas, totalizando 2.700 treinamentos. A repetição com sementes diferentes é importante porque os resultados de redes neurais variam com a inicialização dos pesos, a ordem dos lotes e, no caso das defesas aleatórias, com o próprio sorteio do ruído ou das posições mantidas.

A terceira etapa corresponde ao treinamento dos modelos em um cenário que simula um único cliente, sem agregação entre múltiplos participantes. Esse cenário representa o caso mais favorável ao adversário, pois o gradiente observado corresponde diretamente aos dados de um só cliente. Em cada iteração do treinamento, o gradiente é calculado normalmente, a defesa é aplicada sobre ele e somente o gradiente modificado é utilizado para atualizar os parâmetros. Ao final, registra-se a acurácia de teste de cada modelo e armazenam-se os modelos treinados para a etapa seguinte.

Na quarta etapa, o ataque iDLG é executado sobre um modelo representativo de cada combinação de base de dados, arquitetura e defesa, totalizando 45 modelos. Para cada um deles, um mesmo conjunto de 30 imagens de teste, escolhidas de forma equilibrada entre as classes, é submetido ao ataque, que primeiro infere o rótulo a partir do gradiente e depois otimiza uma imagem artificial até que seu gradiente se aproxime do gradiente original. Para cada tentativa, registram-se o erro quadrático médio entre a imagem original e a reconstruída e se o rótulo inferido estava correto.

A quinta e última etapa reúne as métricas de privacidade e de utilidade de todas as configurações. Cada defesa é comparada ao cenário sem defesa por meio do teste t de Welch, que verifica se a diferença entre as médias é estatisticamente significativa, e do d de Cohen, que indica a magnitude dessa diferença. A partir dessas comparações, interpreta-se o compromisso observado para cada defesa, arquitetura e base de dados.

\section{Organização do Trabalho}

Além desta introdução, o trabalho está organizado em outros quatro capítulos. O Capítulo~2 apresenta a fundamentação teórica necessária para compreender o estudo, abordando o treinamento de redes neurais por gradiente, o aprendizado federado, os ataques de inversão de gradientes, a privacidade diferencial, a compressão de gradientes e as métricas de avaliação, e descreve ainda os trabalhos relacionados e o posicionamento desta pesquisa em relação a eles. O Capítulo~3 detalha o desenvolvimento do trabalho, incluindo o desenho experimental, as bases de dados, as arquiteturas, a configuração de cada mecanismo de defesa, o protocolo de ataque e as métricas empregadas. O Capítulo~4 apresenta e discute os resultados obtidos, com foco no compromisso entre privacidade e utilidade de cada defesa e na influência da base de dados, da arquitetura e do tamanho do lote. Por fim, o Capítulo~5 reúne as considerações finais, as limitações do estudo e as perspectivas de trabalhos futuros.
